\section{Sim-to-Real Transfer}

\subsection{为什么需要仿真}

在真实机器人上做 RL 面临巨大挑战：
\begin{itemize}
    \item 数据收集慢（真实时间运行）
    \item 安全风险（机器人可能损坏自身或环境）
    \item 难以大规模并行化
\end{itemize}

仿真环境解决了这些问题，但引入了\textbf{sim-to-real gap}：仿真和真实世界之间的差异。

\begin{importantbox}{Sim-to-Real Gap 的来源}
\begin{itemize}
    \item \textbf{动力学差异}：仿真器无法完美模拟真实物理（摩擦、接触、变形）
    \item \textbf{感知差异}：仿真的视觉/触觉与真实传感器不同
    \item \textbf{执行器差异}：真实电机有延迟、噪声、非线性特性
    \item \textbf{环境差异}：真实环境有未建模的物体、光照变化等
\end{itemize}
\end{importantbox}

\begin{figure}[H]
\centering
\includegraphics[width=0.9\textwidth]{slides-images/slide-12.jpg}
\caption{讲者用带外参编码器的训练框架说明：在仿真中随机化质量、摩擦、地形等因素，是把策略迁移到真实机器人前的关键准备步骤。}
\end{figure}

\subsection{Domain Randomization}

\begin{importantbox}{域随机化}
Domain Randomization 是弥合 sim-to-real gap 的主要方法之一：

在仿真中随机化各种环境参数（摩擦系数、物体大小、质量、视觉外观、传感器噪声等），使策略在训练阶段就"见过"各种可能的环境变体。

核心假设：如果策略能在足够多样的仿真环境中都表现良好，那么真实世界可以被视为这些变体中的一个。
\end{importantbox}

\begin{warningbox}{Domain Randomization 的局限}
\begin{itemize}
    \item 过度随机化可能导致策略过于保守
    \item 需要确保真实世界参数在随机化范围内
    \item 对某些物理现象（如软体接触），即使随机化也难以覆盖
    \item 选择哪些参数进行随机化以及随机化的范围需要领域知识
\end{itemize}
\end{warningbox}

\subsection{本章小结}

Sim-to-real transfer 是 RL 用于真实机器人的关键桥梁，domain randomization 是目前最常用的方法。

